\section{Circuito local seleccionado para retirada y rearme (CEN-26)}
\label{sec:sd4-circuito-cen26}

Cada camino tiene dos permisos independientes: \texttt{SR-R} para rectificador y \texttt{SR-B} para bancada. Se seleccionan cuatro Schneider \texttt{XPSUAF13AP}, alimentados a 24 V DC; dos Phoenix \texttt{QUINT4-PS/1AC/24DC/10}, referencia 2904601, uno por camino; y cuatro Phoenix \texttt{PTCB E1 24DC/1A NO}, referencia 2909902, uno por módulo. La fuente de control se conecta a la salida protegida de su camino, no detrás de la entrada AC del rectificador que retira. Las salidas positivas A/B no se puentean directamente. ENV-40 incorpora alimentación cruzada sólo para los rieles ambientales mediante dos ORING; permisos y SEL conservan sus ramas originales. Su referencia secundaria común y fallos de retorno se verifican en el unifilar integrado. Se fija 24 V y no se usa la entrada Rem para parar la fuente. Se selecciona protección primaria Schneider \texttt{A9F79210}, iC60N bipolar 10 A curva C, una por fuente. El poder de corte se contrasta con el cortocircuito calculado en ese tablero, sin atribuirle la capacidad del NSX. PTCB admite 18--30 V DC, entrega 1 A y protege el circuito del módulo; su salida se conduce a A1 y el retorno a A2/B2. Su borne IN- es referencia interna, no retorno de carga. El PTCB no es interruptor de batería ni limitador de corriente de carga \parencites[pp.~21--23, 41--42 y 45]{lote26-xpsmanual}[pp.~3--5]{p6-puestos-quint}{lote26-ptcb}{lote26-mcb10}.

XPSUAF dispone de tres contactos de seguridad NO. Se fija función de aplicación 1 y arranque supervisado, posición 3 con dinamización, sin puente Y1--Y2. La parada local usa una cabeza de enclavamiento Schneider \texttt{ZB5AS844} con collar \texttt{ZB5AZ009} y dos bloques \texttt{ZBE102}, de apertura positiva: dos conjuntos por camino, para no compartir los mismos contactos entre dos módulos. Un contacto se conecta S11--S12 y el otro S21--S22. Cada permiso tiene pulsador propio \texttt{XB5AA31} NO entre Y1 e Y2; el pulso debe mantenerse al menos 80 ms y su liberación produce el arranque supervisado. El mando no aplica alimentación a CN12 ni REPO. Z1 sólo entrega diagnóstico al registro y no controla seguridad. La retirada se produce localmente sin WAN/PLC. El proyecto no atribuye PL e/SIL 3 al conjunto por seleccionar un módulo que puede alcanzar esas prestaciones: sensores, cableado, elementos finales y análisis completo deben demostrar la función \parencites[pp.~41--45 y 51--54]{lote26-xpsmanual}{lote26-estop}{lote26-collar}{lote26-contactonc}{lote26-reset}.

\textbf{Rectificador.} El circuito de 230 V de control pasa por fusible gG 10 A Schneider \texttt{DF2CN10} en portafusible \texttt{DF101}, luego por 13--14 y 23--24 de \texttt{SR-R} en serie y termina en la MN \texttt{LV429407} ya seleccionada. El fusible corresponde a la protección externa de contactos prescrita por XPSUAF; no se usa un fusible aM como protección completa del conductor. Un contacto \texttt{29450} del NSX, configurado OF, registra abierto/cerrado: no se lo presenta como contacto espejo de potencia. El disparo de MN requiere rearme manual del interruptor además de rearmar \texttt{SR-R}; una salida recuperada no cierra el NSX. Las entradas de permiso ambiental de banco deben abrir este circuito ante hidrógeno, caudal insuficiente o fallo de supervisión. Su sensor/contacto efectivo y evaluación del fallo soldado se individualizan con ENV; un conector reservado o el guardia de caudal de calefacción no acredita esa entrada. Hasta disponer de esos sensores y contactos no se habilita recarga por una inferencia de presión \parencites[p.~22]{lote26-xpsmanual}{lote26-fusible10}{lote26-fuseholder}{lote26-of}[pp.~C-37 y F-38]{lote23-nsx2026}.

\textbf{Bancada.} Se seleccionan dos Schneider \texttt{C10F3TM016}, NSX100F trifásicos 16 A TM-D, uno por bancada instalada, con MN \texttt{LV429407}. Se añaden dos contactos OF \texttt{29450} para registrar esos interruptores, cuatro OF en total junto con rectificadores. Se conserva la alimentación resistiva trifásica sin neutro de la variante Crestchic, y el PE continuo sin maniobra. Cada salida de interruptor alimenta dos contactores Schneider \texttt{LC1D32P7} en serie, tres polos, bobina 230 V AC/50--60 Hz; cuatro contactores en total. Son elementos de retirada de la carga resistiva. No maniobran el auxiliar monofásico de refrigeración. Las salidas 13--14 y 23--24 de \texttt{SR-B} energizan respectivamente K1 y K2; 33--34 alimenta MN. Se protege la fase común de salidas con \texttt{DF101}/\texttt{DF2CN10}. Cada bobina de contactor requiere 70 VA de captación y 7 VA de mantenimiento a 50 Hz; $70/230=0,3044$ A es inferior a los 3 A AC-15 de cada contacto XPS. Los contactos NC espejo integrados de K1 y K2 se conectan en serie con el pulsador de rearme Y1--Y2: un contactor que no abrió impide rearmar. No se agrega supresor de bobina que altere el tiempo publicado sin recalcularlo \parencites[pp.~1--3]{lote26-contactor}[pp.~1--2]{lote26-nsx16}[p.~22]{lote26-xpsmanual}.

El interruptor de bancada tiene Icu/Ics de 36 kA a 380--415 V y ajuste térmico $0,7$--$1I_n$; se fija $I_r=16$ A. Se proyecta circuito propio de cobre 4G4 mm$^2$, longitud máxima 10 m y cable 0,6/1 kV, con verificación de intensidad admisible corregida, caída, cortocircuito/tiempo y puesta a tierra por el método de instalación definitivo. No se afirma coordinación de tipo 2 NSX/contactores sin tabla específica. La protección térmica no impone instantáneamente un máximo de corriente: la ficha declara 120--400 s a $1,5I_n$ y disparo magnético fijo 190 A. Una selección de 200 kW errónea desde LC10 no queda convertida en un escalón seguro por ese interruptor \parencite[pp.~1--2]{lote26-nsx16}.

\textbf{Tiempo de retirada y separación.} Una solicitud en entrada de XPS desactiva salidas en hasta 20 ms; LC1D32 abre en 4--19 ms. La suma de catálogo de ese recorrido es hasta 39 ms desde entrada, antes del sensor, cables y cualquier lógica previa. La MN publica 50 ms de respuesta del accesorio, no tiempo integral de corte del interruptor. Ante pérdida de 24 V del módulo, el manual XPSUAF1 declara hasta 120 ms, no 20 ms: el camino a contactor resulta hasta 139 ms desde esa pérdida. Tampoco se suma un sostenimiento típico de QUINT como si fuera cota máxima. La aceptación mide el recorrido completo de cada causa hasta apertura y verifica el estado espejo. El auxiliar de bancada se alimenta por circuito independiente, permanece durante enfriamiento y no sigue K1/K2/MN. Extracción de batería y detección permanecen en sus circuitos protegidos independientes. Retirar \texttt{SR-B} no retira \texttt{SR-R}; retirar rectificador consume batería y requiere continuidad de cargas y camino sano dentro de su reserva \parencites[p.~23]{lote26-xpsmanual}[p.~2]{lote26-contactor}[p.~C-37]{lote23-nsx2026}.

\section{Medición de grupo y frontera del desarrollo GCI-26}
\label{sec:sd4-medicion-gci26}
GCI-29 sustituye los seis RM35 por dos SEL-751, conservando seis TC y circuitos de retirada de CEN-26. La imprecisión de RM35 y su inhibición de arranque motivan el cambio; no se adquieren ni se conservan en paralelo como permiso adicional. El nuevo desarrollo fija umbrales por fase y tensión, errores y latencias, sin declarar límite instantáneo ni configuración ejecutable ya validada. El circuito de 16 A mantiene solicitud manual máxima de diez pasos, sujeta al total del grupo y fase; dos LC1D32P7 por bancada y enfriamiento independiente siguen vigentes. Las entradas de gas y caudal útil no se sustituyen por la medida eléctrica.

\section{Recarga y demanda de control reconciliadas (ECI-26)}
\label{sec:sd4-recarga-eci26}

La guía 93T identifica el máximo de 60 A como configurable, pero no publica en esa tabla la resolución, límites inferiores ni una asignación autorizada de inhibición selectiva a CN12. Se mantiene como ajuste propio requerido 20 A agregados para estudio de recuperación; no se declara un perfil de fábrica ya autorizado. Con una sola cadena disponible, aun el ajuste requerido de 20 A sería menor que los 54 A de CSB, pero sigue necesitando supervisión de corriente/tensión/temperatura y aceptación del cargador real. La MN se reserva a emergencia y no se cicla para sostener ese ajuste; tampoco resuelve una cadena con corriente anormal mientras la otra oculta el total. A 588 V/20 A se obtiene techo de 11,76 kW DC. Para reponer 63,8298 kWh retirados en el modelo de 60 kW/60 min, el mínimo ideal es $63,8298/11,76=5,4277$ h, sin absorción ni pérdidas; a 30 min sería 2,7139 h. Esas cotas no son plazo de recarga ofrecido y veinte amperios pueden exceder la entrada permitida cuando el servicio está a plena carga. Se conserva extracción dimensionada para la envolvente precedente; no se reduce por usar un ajuste aún no autorizado \parencites[p.~98]{p7-eaton-93t-guia}{csbHRL12540WRA260131}.

La referencia DC anterior, sustituida por 72,8 W antes de sensores en ECI-29, es ocho vatios para cuatro XPS, hasta 3,2 W de cuatro PTCB y 3,6 W de seis RM35: 14,8 W, distribuida por camino a 7,4 W antes de sensores y cableado. Se reservan, sin convertirlos en consumo real a carga parcial, 274 VA de entrada nominal por QUINT: 548 VA entre dos fuentes. En régimen de bancada retirada, sólo las dos MN de rectificador añaden 10 VA continuos, total de reserva 558 VA. Con dos bancadas permitidas, cuatro contactores añaden 28 VA y cuatro MN suman 20 VA: reserva 596 VA. La captación de cada bancada añade dos veces 70 VA y la cota de cableado MN de 30 VA, 170 VA por rama; al rearmar también rectificador se reservan 200 VA por rama, 400 VA entre ambas, además de las fuentes. No se suman los 14,8 W DC otra vez al lado AC ni se toma 93,4\,\% típico de QUINT como eficiencia a 7,4 W. Estas dos fuentes son partidas nuevas de seguridad, diferentes de las QUINT de puestos; las bobinas MN precedentes quedan reemplazadas por este desglose, no duplicadas \parencites[pp.~3--4]{p6-puestos-quint}{lote26-ptcb}[p.~21]{lote26-xpsmanual}[p.~1]{lote26-rm35}[p.~2]{lote26-contactor}[pp.~C-37 y E-20]{lote23-nsx2026}.

El tratamiento de ENV-25 selecciona dos recalentadores de 9 kW activos y reserva 100 W de disipación de dos TTC25: 18,100 kW de tratamiento, cuatro unidades instaladas con dos activas. Sobre los 50,360015 kW parciales anteriores resulta 68,460015 kW sin margen adicional del tratamiento, o $50,360015+1,20(18,100)=72,080015$ kW con ese margen. La reserva continua máxima de 0,596 kVA de este control se distribuye en el balance por camino y estado, sin sumarla como 0,596 kW medidos ni restarla de TI. Frío, bombas, ventiladores de unidades y receptores de entrada única siguen requiriendo selección integrada. No se libera una UPS mayor por estas sumas parciales ni se ofrece PUE completo desde cotas de compra. La recepción incluye descarga--recarga--nuevo corte y clima completo, además de retirada por cada causa, estado del camino sano, corriente por cadena y fase, tiempos y energía separados.
